\documentclass[conference]{IEEEtran}
\IEEEoverridecommandlockouts
% The preceding line is only needed to identify funding in the first footnote. If that is unneeded, please comment it out.
\usepackage{cite}
\usepackage{amsmath,amssymb,amsfonts}
\usepackage{algorithmic}
\usepackage{graphicx}
\usepackage{textcomp}
\usepackage{xcolor}
\def\BibTeX{{\rm B\kern-.05em{\sc i\kern-.025em b}\kern-.08em
    T\kern-.1667em\lower.7ex\hbox{E}\kern-.125emX}}
\begin{document}

\title{Ironbound Tactics\\
%{\footnotesize \textsuperscript{*}Note: Sub-titles are not captured in Xplore and
%should not be used}

}

\author{\IEEEauthorblockN{\textsuperscript{} Julian Kalb}
\IEEEauthorblockA{\textit{Game Lab III} \\
\textit{Julius-Maximilians-Universität}\\
Würzburg, Germany}

}

\maketitle

\begin{abstract}
For the scope of the  GameLab III module, which is part of the bachelors degree Games Engineering, the students
were tasked to develop a video game in the programming language C++, without resorting to Game Engines
or engine-like libraries.
The resulting game is called Ironbound Tactics.
It is a turn based strategy game, where the player fights with their own composed army against randomly selected
armies of a computer-controlled opponent to set new highscores and unlock rewards.
 
\end{abstract}

%\begin{IEEEkeywords}
%component, formatting, style, styling, insert
%\end{IEEEkeywords}

\section{Introduction}
Modern game-engines offer a variety of functionality and features, which simplify  the development of video games
considerably.
They come with a toolbox, that by default takes on tasks like rendering, physics, asset-management and
input-handling, which gives developers more freedom to focus on other aspects like gameplay, level- and
character design.
For this project on the other hand, the challenge was to renounce those features, and program a game from scratch.
During a workshop, where all students were guided to find a fitting concept for the game, the decision was made to create a
turn-based, "roguelike" game, where the player would, at the beginning of each run, create their own army and then proceed
to fight randomly selected opponents.
The idea was to offer the player a set of unit-types, which differ in attack speed, accuracy, damage and health points.
The Player would then fight randomly picked armies of computer-controlled opponents, until they were defeated.
The general inspiration for Ironbound Tactics, especially the turn-based combat, stemmed from the Pokémon \cite{b1}
franchise, though far less complicated.

\section{Conception Phase}
The decision for Ironbound Tactics to become a turn-based strategy game already came with some design implications.
The gameplay of the final game would be based on existing entries of the "JRPG" genre.
The core gameplay loop would work as follows: The turn of the player starts, the player decides what actions to take,
choose the according options in the user interface and end their turn.
Then damage calculation and attack animations would be executed.
After that the application would evaluate, whether player or opponent had lost.
Finally, either the current fight would end, or the next turn would start, beginning the cycle anew.
Thus,the user interface seemed like the logical entry point to the project.

\subsection{User Interface}
The decision was made to use an existing library for the Graphical User Interface (GUI), since the custom
implementation was deemed too time-consuming.
The first Prototype was created using the library imgui \cite{b2}.
The library was easy to set up and use, which was important to get started and test the first equation for damage
calculation.
And while imgui~\cite{b2} was a good choice for prototyping, there was one big incompatibility: The ability to skin
the GUI objects of imgui~\cite{b2} was very limited.
Therefore, at the end of the prototyping phase the decision was made to switch to another library called Nuklear~\cite{b3}.

\subsection{Gameplay}
The general idea of the gameplay was to let the player create their own constellation of units, by letting them choose
out of a predefined pool of units, creating their army.
There would be some kind of limitation for how many units the player would be allowed to choose, both per individual
unit and in total.
Three different types of units were implemented in the final version of Ironbound Tactics:
\begin{itemize}
    \item \textbf{Archer} units have the lowest accuracy and individual damage but deal damage at the end of the turn.
    \item \textbf{Infantry} units are melee units.
    They deal higher damage than \textbf{Archers}, but need to walk to their target and therefore need two turns to
    deal damage.
    \item \textbf{Catapult} units are the slowest, dealing damage at the end of the third turn after declaring an attack.
    This unit also has a cost to be used, making it necessary to first assigning a fix amount of either \textbf{Archers}
    or \textbf{Infantry} units to use them.
    To balance these drawbacks \textbf{Catapults} deal the highest damage per instance.
\end{itemize}

When the player is done with the distribution of their units, the first round of the combat loop would start.
During a round of combat the player would fight the army of a computer-controlled opponent.
This army would be randomly selected out of a set pool of presets.
Both armies would fight, until all units of either side were exhausted.
If the players army is the first to be exhausted, the game would end and quit.
On the other hand, if the opponents army is the first to be exhausted, the player would get the chance to select
a reward, before entering the next round against a new opponent with a new army.

\subsection{Style}
During the creation of the first prototype the decision for the art style and visual setting were made.
The art style would be 2D Pixelart, since fitting assets were available from the start.
The setting would be medieval.
This decision was heavily implicated by the general game concept, since the mixture of close- and ranged-combat
seemed to offer good variety during unit design.

\section{System Design}
Sources and instructions to implement a rendering pipeline and create custom systems were provided to the students.
Since the game was to be created without the usage of an engine, the design of those systems was fairly complex.

\subsection{Entity Component System}
An entity component system (ECS) is a pattern, that is commonly used in game development.
Commonly object-oriented programming and inheritance are used as an architectural pattern during software development.
In object orientation entities are designed similarly to the perceived world, where base classes are created with
attributes, always considering reusability in similar classes.
An entity would therefore be seen as an object, that is an instance of such a class.
In an entity component system on the other hand, entities are typically represented by lightweight identifiers.
This design philosophy entails separating entities, components and systems.
Components are assigned to entities, which are identified by a unique identifier.
The assigned components are then stored into dedicated containers for their type.
This is advantageous in settings like game development, since often entities have a lot of attributes in common, like
a mesh and a texture, and only differ in small items.
As those common attributes can often be processed at the same point in the program, the resulting serialization of
components into dedicated containers can give a significant boost in performance.

A basic custom implementation of an ECS was provided to the students, which was expanded during the development
of Ironbound Tactics.
While a lot of created systems are only useful to this game, some concepts like the sound-system and the render-system
were designed to be reusable and applicable to other types of games.

\subsection{Rendering}
In the instructions OpenGL~\cite{b4} and GLFW~\cite{b5} were used to set up fundamental functionality like creating
a window and a basic render pipeline.
As Ironbound Tactics was planned as a 2D game, the complexity and workload of the render pipeline would
be on the lower end.
Therefore it was decided to keep using OpenGL~\cite{b4}.
Still there were some adjustments to be made.

As an example, the starting point was making a draw call to the graphics card for each individual entity.
Since Ironbound Tactics would be likely to have a lot of instances of the same unit be at the screen at the same time,
this would become a problem, as individual draw-calls are a very demanding process on the computer.
A component "Instance Buffer" was therefore introduced, which, in combination with a scene-graph structure, allowed to
batch duplicate instances of a single entity, into a single buffer, reducing the necessary draw-calls down to one.

\section{Testing \& Feedback}
During the course of the GameLab III module, the students had the opportunity to attend to regular feedback sessions.
From the start this feedback was very helpfully in shaping and refining Ironbound Tactics.
For example, for the longest time the game only consisted of the user interface, as the logic for turn handling and
damage calculation was a difficult task.
For this reason development did get stuck in this phase for a while.
During one of the feedback sessions the advice was given, to not postpone the visual representation too far to the end
of the development, as the visual feedback would help getting a grasp of the impact of what was going on behind the
scenes.
Implementing this feedback provided another point to continue to work from, resulting in a decrease of frustration.\newline

Ironbound Tactics was also tested by peers.
The received feedback did not only result in improved balancing, the discovery, reproduction and fix of bugs, but also in refining existing
features and coming up with new ones.
As an example, it was noted, that the available decisions the player has during their turn were to ambiguous.
The idea to create icons, that hover over the possible actions was given as a possible solution, which was well received
after implementation.
Also, the replayability was criticised, as the motivation to chase a new highscore and experiment with different setups
and rewards died down quickly.
A tester came up with the idea, to give the player the ability to unlock the textures of your opponents army.
Since this feedback came at a very late stage of development it had to be evaluated, if the implementation of this
feature would still be possible.
It turned out to be not too complex, as the implemented systems were simple to extend.

\section{Conclusion}
Ironbound Tactics was, in its final state, again reviewed by peers.
The general tone of the feedback was very positive.
Players enjoyed the art style and general gameplay loop, although there was still some criticism of the balancing.
Also, some minor bugs persisted until the final release.

Concluding the end of the GameLab III module, Ironbound Tactics can be described as a fun experience.
In the future the game could be expanded in several ways.
For example an idea to update the graphics style from 2D to "2.5D" was also scrapped during development.
This would improve the perceived visual quality of the game a lot, and also open up the possibility to add more lighting
effects.
Also, it would be possible to expand the existing selection of units with broader variety, making the
gameplay experience more diverse each run.

\begin{thebibliography}{00}
\bibitem{b1} Game Freak, Pokémon Red Version, Game Boy. Nintendo, 1996.
\bibitem{b2} O. Ocornut, “Dear ImGui”, GitHub repository. [Online]. Available: https://github.com/ocornut/imgui
\bibitem{b3} V.D. (vurtun), "Nuklear - Immediate-Mode-UI", GitHub repository. [Online]. Available: https://github.com/Immediate-Mode-UI/Nuklear
\bibitem{b4} Khronos Group, OpenGL 4.6 Core Profile Specification, Khronos Group, 2017.
\bibitem{b5} M. Geelnard and C. Heine, “GLFW — An OpenGL framework,” GLFW
\end{thebibliography}
\vspace{12pt}
\end{document}
